\section{Determinant mass and the sharp lower-spectrum conversion}
\label{sec:tpc67-determinant-mass}

The preceding interface fixes the literal normalized full-output matrix
\(B_X\), its unit-diagonal Gram
\[
 K_X=B_X^*B_X,
\]
the number \(r_X\ge1\) of active physical rows, and the full determinant
mass
\[
 \mathfrak D_X
 =
 \sum_{\substack{I\subseteq\mathcal Y_X^{\mathrm{full}}\\|I|=r_X}}
 |\det B_X[I]|^2.
\]
No output quotient or selected-row restriction is made in this section.

\subsection{The exact Cauchy--Binet identity}

\begin{theorem}[Full determinant mass is the Gram determinant]
\label{thm:tpc67-cauchy-binet-mass}
For every \(r_X\ge1\),
\begin{equation}
 \boxed{\qquad
 \mathfrak D_X=\det(B_X^*B_X)=\det K_X.
 \qquad}
 \label{eq:tpc67-cauchy-binet-mass}
\end{equation}
The identity also holds when
\(|\mathcal Y_X^{\mathrm{full}}|<r_X\): both sides then vanish.
\end{theorem}

\begin{proof}
Fix the orderings from the full-minor interface.  Cauchy--Binet applied
to \(B_X^*B_X\) gives
\[
 \det(B_X^*B_X)
 =
 \sum_{\substack{I\subseteq\mathcal Y_X^{\mathrm{full}}\\|I|=r_X}}
 \det\!\bigl(B_X[I]^*\bigr)\det B_X[I].
\]
Each summand is \(|\det B_X[I]|^2\), which proves
\eqref{eq:tpc67-cauchy-binet-mass}.  If the output space has dimension
less than \(r_X\), then \(B_X\) has rank less than \(r_X\), while the
minor sum is empty by convention; hence both sides are zero.
\end{proof}

For later use, define
\begin{equation}
 c_1:=1,
 \qquad
 c_r:=\left(\frac{r-1}{r}\right)^{r-1}
 \quad(r\ge2).
 \label{eq:tpc67-sharp-determinant-constant}
\end{equation}

\begin{theorem}[Sharp determinant-to-gap conversion]
\label{thm:tpc67-sharp-determinant-gap}
Let \(K\succeq0\) be an \(r\times r\) matrix with unit diagonal, and put
\(\beta=\lambda_{\min}(K)\).  If \(r=1\), then
\[
 \beta=\det K=1.
\]
If \(r\ge2\), then
\begin{equation}
 \det K
 \le
 \beta
 \left(\frac{r-\beta}{r-1}\right)^{r-1}
 \le
 \frac{\beta}{c_r}.
 \label{eq:tpc67-refined-determinant-upper}
\end{equation}
Consequently, for the literal matrix,
\begin{equation}
 \boxed{\qquad
 \beta_X
 \ge
 c_{r_X}\mathfrak D_X
 \ge
 e^{-1}\mathfrak D_X.
 \qquad}
 \label{eq:tpc67-determinant-gap-lower}
\end{equation}
For every fixed \(r\ge2\), \(c_r\) is the largest universal constant
for which
\[
 \lambda_{\min}(K)\ge c_r\det K
\]
holds for every unit-diagonal positive semidefinite \(K\).
\end{theorem}

\begin{proof}
The case \(r=1\) is immediate.  Suppose \(r\ge2\), and list the
eigenvalues as
\[
 0\le\lambda_1=\beta\le\lambda_2\le\cdots\le\lambda_r.
\]
The unit diagonal gives
\[
 \sum_{j=1}^r\lambda_j=\operatorname{tr}K=r.
\]
The arithmetic--geometric mean inequality applied to
\(\lambda_2,\ldots,\lambda_r\) yields
\[
 \prod_{j=2}^r\lambda_j
 \le
 \left(\frac{r-\beta}{r-1}\right)^{r-1}
 \le
 \left(\frac r{r-1}\right)^{r-1}.
\]
Multiplication by \(\beta\) proves
\eqref{eq:tpc67-refined-determinant-upper}.  Combine it with
\cref{thm:tpc67-cauchy-binet-mass} to obtain the first inequality in
\eqref{eq:tpc67-determinant-gap-lower}.  Moreover,
\[
 c_r
 =
 \left(1+\frac1{r-1}\right)^{-(r-1)}
 \ge e^{-1},
\]
which proves the second.

It remains to prove sharpness.  For \(0<\varepsilon\le1\), let
\begin{equation}
 K_{r,\varepsilon}
 :=
 \frac{r-\varepsilon}{r-1}I_r
 -
 \frac{1-\varepsilon}{r-1}\mathbf1\mathbf1^*.
 \label{eq:tpc67-negative-equicorrelation}
\end{equation}
This is the negative equicorrelation matrix with diagonal one.  Its
eigenvalue on \(\mathbf1\) is \(\varepsilon\), and every vector
orthogonal to \(\mathbf1\) has eigenvalue
\((r-\varepsilon)/(r-1)\).  Thus it is positive definite and
\begin{align*}
 \lambda_{\min}(K_{r,\varepsilon})
 &=\varepsilon,\\
 \det K_{r,\varepsilon}
 &=
 \varepsilon
 \left(\frac{r-\varepsilon}{r-1}\right)^{r-1}.
\end{align*}
Therefore
\[
 \frac{\lambda_{\min}(K_{r,\varepsilon})}
      {\det K_{r,\varepsilon}}
 =
 \left(\frac{r-1}{r-\varepsilon}\right)^{r-1}
 \longrightarrow c_r
 \qquad(\varepsilon\downarrow0).
\]
Any universal multiplier larger than \(c_r\) fails for sufficiently
small \(\varepsilon\).
\end{proof}

\begin{remark}[What sharpness does and does not say]
\label{rem:tpc67-determinant-sharpness-scope}
The negative equicorrelation family proves sharpness among all
unit-diagonal Gram matrices of a fixed dimension.  It does not assert
that the literal fixed-\(h_0\) matrix has this form.  Any improvement for
the physical carrier must use additional actual-support or coefficient
structure.
\end{remark}

\subsection{Summing certified full minors}

\begin{corollary}[Certified-minor mass]
\label{cor:tpc67-certified-minor-sum}
Let
\[
 \mathscr I_X^{\mathrm{cert}}
 \subseteq
 \{I\subseteq\mathcal Y_X^{\mathrm{full}}:|I|=r_X\},
\]
and suppose that for each \(I\in\mathscr I_X^{\mathrm{cert}}\) an actual
coefficient calculation supplies a number \(\eta_X(I)\ge0\) such that
\begin{equation}
 |\det B_X[I]|\ge\eta_X(I).
 \label{eq:tpc67-certified-one-minor}
\end{equation}
Then
\begin{equation}
 \boxed{\qquad
 \mathfrak D_X
 \ge
 \sum_{I\in\mathscr I_X^{\mathrm{cert}}}\eta_X(I)^2,
 \qquad
 \beta_X
 \ge
 c_{r_X}
 \sum_{I\in\mathscr I_X^{\mathrm{cert}}}\eta_X(I)^2.
 \qquad}
 \label{eq:tpc67-certified-minor-mass}
\end{equation}
In particular, if \(L_X\) distinct certified minors all satisfy
\(|\det B_X[I]|\ge\eta_X\), then
\begin{equation}
 \beta_X\ge c_{r_X}L_X\eta_X^2.
 \label{eq:tpc67-uniform-certified-minors}
\end{equation}
\end{corollary}

\begin{proof}
Retain only the indicated nonnegative summands in
\eqref{eq:tpc67-cauchy-binet-mass}, apply
\eqref{eq:tpc67-certified-one-minor}, and then use
\cref{thm:tpc67-sharp-determinant-gap}.
\end{proof}

\begin{remark}[Activity labels are not numerical certificates]
\label{rem:tpc67-activity-needs-minor-bound}
The preceding forced-matching and cycle-activity analysis can be used to
construct \(\mathscr I_X^{\mathrm{cert}}\).  For example, if the minor
\(B_X[I]\), a matching \(\phi\), and its forced factors satisfy
\cref{thm:tpc67-activity-certificate}, then one valid choice is
\begin{equation}
 \eta_X(I)
 :=
 P_{\mathrm{for}}(I)P_\phi(I)
 \bigl(1-\mathcal A_\phi(I)\bigr)_+.
 \label{eq:tpc67-activity-certified-eta}
\end{equation}
However, a combinatorial activity label by itself is not an admissible
\(\eta_X(I)\).  The corollary begins only after the literal determinant
expansion has been shown to have the numerical lower bound
\eqref{eq:tpc67-certified-one-minor}.  Conversely, minors not covered by
the chosen certificate family remain in \(\mathfrak D_X\); omitting them
from the certified sum cannot be used as evidence that they vanish.
\end{remark}

\begin{remark}[Claim level]
\label{rem:tpc67-determinant-mass-claim-level}
\Cref{thm:tpc67-cauchy-binet-mass,thm:tpc67-sharp-determinant-gap,cor:tpc67-certified-minor-sum}
are finite-dimensional L0 statements.
Their L1 specialization uses the already fixed literal full-output
matrix \(B_X\).  They do not prove that a certified full minor exists,
that the certified mass is polynomially large, or that
\(\mathfrak D_X\) avoids exponential decay as \(r_X\) grows.  No parity
or prime-pair conclusion follows from the determinant identity alone.
\end{remark}
